% ch09 — Automatic Differentiation
\chapter{Automatic Differentiation}
\label{ch:autodiff}

Computing Christoffel symbols requires derivatives of the metric components.
Rather than hand-deriving and hard-coding these derivatives for each spacetime,
the codebase uses forward-mode automatic differentiation (AD) with dual
numbers. This chapter explains the technique, compares it with finite
differences and analytic expressions, and briefly surveys reverse-mode AD.

\section{Forward-Mode AD and Dual Numbers}
\label{sec:ad-forward}
\coderef{Spacetime}{AutoDiff}

We define the dual-number algebra $a + b\,\varepsilon$ with
$\varepsilon^2 = 0$, show how arithmetic operations propagate derivatives
exactly, and demonstrate the implementation in the codebase's \texttt{AutoDiff}
module.

\section{AD for Christoffel Symbols}
\label{sec:ad-christoffels}

We show how seeding the metric function with dual numbers automatically
produces the metric derivatives needed to compute Christoffel symbols,
eliminating the need for analytic derivative expressions.

\section{Comparison: AD, Finite Differences, and Analytic Derivatives}
\label{sec:ad-comparison}

We compare the three approaches in terms of accuracy, computational cost, and
implementation effort, demonstrating that AD achieves machine-precision
derivatives at a modest overhead.

\section{Reverse-Mode AD\starmark{}}
\label{sec:ad-reverse}

We outline reverse-mode (adjoint) AD for completeness, noting its advantages
for functions with many inputs and few outputs, and discuss potential
applications in gradient-based optimisation of spacetime parameters.

\paragraph{Key References.}
The dual-number approach follows \cite{Griewank2008}; the application to
GR is discussed in \cite{Liska2018}; reverse-mode AD is surveyed in
\cite{Baydin2018}.
